\chapter{理论基础与研究环境}
\label{chap:foundations}

本文以空间公共物品博弈作为合作演化环境，通过个体强化学习描述策略的逐步调整，并研究邻居信息进入学习过程的方式。本章介绍后续模型共同使用的概念与评价口径。需要区分三个层次：博弈规则决定真实收益，学习规则决定个体如何调整策略，消息机制决定个体能够获得怎样的邻居信息。将三个层次分别定义，有助于识别合作变化究竟来自环境激励、学习过程还是观测质量。

\section{空间公共物品博弈与合作演化}
\subsection{个体成本与群体收益}

公共物品博弈描述一种共享收益与个体成本并存的决策情形。设一个群组包含 $k$ 个成员，其中有 $n_C$ 个成员合作，每个合作者投入成本 $c$，公共投入经协同因子 $r$ 放大后均分。对于已经确定的群体动作配置，合作与背叛成员分别获得
\begin{equation}
 \pi_C=\frac{rcn_C}{k}-c,\qquad
 \pi_D=\frac{rcn_C}{k}.
 \label{eq:foundation-group-payoff}
\end{equation}
这两个量描述同一组内不同成员的所得。在讨论某个成员改变动作的激励时，还需要同时改变 $n_C$，不能把固定 $n_C$ 下的差值直接当作单边改变动作的收益差。

固定其余成员的动作，一个成员由背叛改为合作，使公共投入增加 $c$，自身增加的分红为 $rc/k$，同时付出成本 $c$。其净变化为 $c(r/k-1)$。当 $1<r<k$ 时，该变化为负，而整个群组的总收益增加 $c(r-1)>0$。这一计算说明了个体激励与群体收益之间的冲突：合作有利于群体，但仅考虑一次交互的自身收益时，个体可能选择背叛。

空间公共物品博弈把上述交互嵌入网络。本文沿用工作1的周期方格环境\cite{kang2025neighbor}：每个节点具有四个一阶邻居，每个组由中心节点及其一阶邻居组成，智能体参与五个重叠群组。个体收益按所参与的群组累积，因此某个位置的收益不仅与本节点动作有关，也与周围多个群组的合作结构有关。周期边界消除网格边缘的度数差异，使各节点具有一致的交互机会。

\subsection{空间分布与信息范围}

合作率描述合作者占比，空间分布描述这些合作者位于何处。相同合作率可以对应分散的合作个体，也可以对应连续的合作区域，两者所处的局部收益条件未必相同。因此，本文在考察群体合作水平的同时，也使用策略或声誉快照观察空间组织。

交互范围和信息范围具有不同含义。交互范围决定公共物品组的组成以及成本、收益如何计算；信息范围决定智能体能观察哪些节点。本文的感知参数 $M$ 改变后一种范围。在 $M=2$ 时，一个智能体可获得十二个邻居的信息，但每个博弈组仍然只有五个成员。扩大信息范围本身没有增加公共投入，也没有改变每组收益分配规则。

\section{强化学习与动作价值更新}
\subsection{综合奖励与合作的即时激励}

个体学习所使用的奖励不仅影响数值尺度，也可能改变动作之间的即时激励。固定其他节点的动作，一个节点由背叛改为合作，会同时改变其参与的五个群组。每组内自身分红增加$r/5$、成本增加1，因此原始累计收益的变化为$r-5$。采用第三章的归一化和声誉增量奖励，综合奖励之差为
\begin{equation}
 R_i(C)-R_i(D)=w_P\frac{r-5}{3r+5}+(1-w_P).
 \label{eq:foundation-marginal-reward}
\end{equation}
该式比较同一轮、相同其他节点动作下的两种反事实选择；没有把之后的声誉状态和邻居反应固定为真实长期轨迹。在$1<r<5$时，纯收益奖励$w_P=1$保留了合作的当期成本劣势。加入声誉增量奖励会减弱这一劣势，当$w_P\leq(3r+5)/(2r+10)$时，还可使合作的即时综合奖励不低于背叛。

例如，$r=3.0$时的临界权重为0.875；工作1采用$w_P=0.95$的部分实验仍处于即时合作奖励较低的情形。因而该条件下的长期合作不能仅用一次动作的即时优势解释。工作2主实验取$w_P=1$，使新增信息控制在保留原始个体成本冲突的环境中接受检验。式\eqref{eq:foundation-marginal-reward}还说明，声誉奖励不应仅被称作无影响的数值变换，其行为作用需与邻居附加修正分别讨论。

\subsection{状态、动作、奖励与策略}

强化学习关注智能体通过交互反馈调整行为的过程\cite{sutton2018reinforcement}。在离散时刻 $t$，智能体获得状态表示 $s_t$，依照策略 $\pi(a\mid s_t)$ 选择动作 $a_t$，随后获得奖励并进入下一状态。动作价值 $Q^\pi(s,a)$ 表示在状态 $s$ 下先执行动作 $a$、之后遵循策略 $\pi$ 时的预期折扣回报。采用本文的当轮奖励记号，可写为
\begin{equation}
 Q^\pi(s,a)=\mathbb E_\pi\left[\sum_{\tau=0}^{\infty}\gamma^\tau R_{t+\tau}
             \mid s_t=s,a_t=a\right],\quad 0\leq\gamma<1.
 \label{eq:foundation-return}
\end{equation}
折扣因子 $\gamma$ 控制后续奖励的权重。即时奖励是一次交互的反馈，动作价值则汇总对未来回报的估计，两者不能互换。

本文进一步区分博弈原始收益 $P_i(t)$、个体学习奖励 $R_i(t)$ 和控制器训练目标 $J(\theta)$。原始收益由公共物品博弈产生；学习奖励可以包括收益归一化与声誉项；控制器训练目标则评价完整多智能体轨迹。三者可能相关，甚至在特定条件下成比例，但定义与使用环节不同，必须分别说明。

\subsection{Q-learning与探索}

Q-learning利用当前经验和下一状态的价值估计更新动作价值\cite{watkins1992q}。对于智能体 $i$ 的一次转移，其时序差分目标为 $R_i(t)+\gamma\max_{a'}Q_i(s_i(t+1),a')$，更新为
\begin{equation}
 \begin{aligned}
 Q_i(s_i(t),a_i(t))\leftarrow Q_i(s_i(t),a_i(t))
 +\eta\bigl[&R_i(t)+\gamma\max_{a'}Q_i(s_i(t+1),a')\\
             &-Q_i(s_i(t),a_i(t))\bigr].
 \end{aligned}
 \label{eq:foundation-q}
\end{equation}
其中，$\eta$ 为学习率。本文的个体模型采用小规模离散 Q 表，便于直接追踪社会信息造成的附加价值变化，而不需要把动作决策交给大型神经网络。

如果始终选择当前估计最优的动作，智能体可能缺少对其他动作的经验。$\epsilon$-greedy 策略用一定概率进行随机探索，其余情况下选择 Q 值最大的动作。后续模型采用随时间下降但保留正下限的探索率。因此，即使群体已经形成较稳定的合作比例，个别节点仍可能因为探索而切换动作。某一时刻出现全合作或全背叛，也不必然构成吸收状态。

\subsection{多个智能体同时学习的限制}

多智能体强化学习的任务分类取决于各个体的目标关系，信息结构则取决于训练和执行阶段能够访问的观测\cite{zhang2021multiagent}。本文虽然以促进合作为研究目标，但个体Q表接收的是各自奖励，不能据此把博弈视为所有节点获得相同回报的完全合作任务。共享控制器通过群体福利选参，个体仍按自身经验在线学习，这两个优化层次并不相同。

经典 Q-learning 的收敛分析建立在相应受控马尔可夫环境、充分访问及学习率条件上\cite{watkins1992q}。本文中的智能体同时更新策略，从单个智能体的角度看，邻居行为分布会随学习而变化。此外，邻域平均声誉的二值表示压缩了大量空间与历史信息，不同真实环境可以映射到同一状态。因此，不能仅因为采用了 Q-learning 更新式，就把经典结论直接用于整个合作演化系统。

本文据此把有限时间内的合作轨迹和策略切换作为经验观测，避免将曲线变平等同于已证明的最优策略收敛。邻居附加更新改变了标准价值迭代过程，也需要在其具体定义与实验范围内评价。第三章分析已发表的固定规则，第四章进一步引入可训练的信息控制器。

\section{社会信息与观测边界}
\subsection{真实行为与报告消息}

设智能体 $j$ 的真实动作为 $a_j(t)$、真实综合奖励为 $R_j(t)$，发送给邻居的消息为
\begin{equation}
 \widehat m_j(t)=\bigl(\widehat a_j(t),\widehat R_j(t)\bigr).
 \label{eq:foundation-message}
\end{equation}
可靠条件下，报告值与真实值一致。消息失真时，接收者获得的动作或奖励可能不同于真实值。区分二者的目的，是使攻击作用在明确的信息通道上，而不同时改变物理博弈环境和个体自身奖励。

通信鲁棒性研究已经考虑了评估消息可靠性并降低有害消息影响的思路。例如，ADMAC通过局部观测与历史信息估计消息的可靠程度，再调节消息对动作偏好的贡献\cite{yu2024admac}。本文关注的具体过程有所不同：智能体继续在线更新自身 Q 表，邻居消息用于附加价值修正。后续实验需要观察错误信息如何积累为策略变化，而不仅是一次动作输出是否改变。

\subsection{本研究的受扰范围}

工作2的基础设置仅干扰 NI 消息通道。异常发送者仍执行与其他节点相同的物理博弈和学习规则，但按照指定规则修改对外发送的动作、奖励报告。自身真实收益、Q 表以及状态所用的声誉汇总不被该攻击直接修改。声誉汇总在这一设置中作为独立可信观测提供，这是方法评价的明确范围假设。

该假设使实验能够集中检验“错误邻居消息进入价值修正”这一作用路径，但也限制了结论。模型不能据此宣称抵抗任意传感器失真、Q 表篡改或全部声誉信息污染。执行时的信息选择模块只读取接收者自身经验和收到的消息；异常节点身份及邻居真实值只供仿真生成攻击与计算评价指标使用。

\section{评价指标与统计口径}
\subsection{合作水平与真实福利}

设网络共有 $N$ 个节点，时刻 $t$ 的合作率为 $f_C(t)=N^{-1}\sum_i\mathbb I\{a_i(t)=C\}$。对于固定评价窗口 $\mathcal T$，平均合作率定义为
\begin{equation}
 \bar f_C(\mathcal T)=\frac{1}{|\mathcal T|}\sum_{t\in\mathcal T}f_C(t).
 \label{eq:foundation-mean-coop}
\end{equation}
全程平均刻画整个学习过程的表现，末段平均刻画指定后期窗口的表现。两者回答不同问题：一个方法可能恢复较早但最终水平较低，另一个方法可能前期低迷但后期达到较高合作。因此，窗口必须预先定义并在方法之间保持一致。

真实群体福利定义为 $W(t)=\sum_iP_i(t)$，不使用可能被篡改的报告奖励。对于本文的五人重叠群组和每组合作成本1，每个组的总收益为 $(r-1)n_C(g,t)$，每个合作者恰好参与五个组。因此
\begin{equation}
 W(t)=5(r-1)N_C(t),\qquad \overline P(t)=\frac{W(t)}{N}=5(r-1)f_C(t).
 \label{eq:foundation-welfare-identity}
\end{equation}
这是由环境定义得到的恒等式。固定 $r>1$ 和相同时间窗口时，平均真实福利与平均合作率成正比。后续训练若采用按 $5(r-1)$ 标准化的福利，其数值就等于相应合作率；不应仅因目标函数使用“福利”一词，就将其解释为与合作水平无关的目标。跨 $r$ 汇总未经标准化的福利则会赋予不同环境不同权重。

\subsection{策略切换与信息使用}

策略切换率衡量相邻时刻发生动作变化的节点比例：
\begin{equation}
 f_{\mathrm{sw}}(t)=\frac{1}{N}\sum_i\mathbb I\{a_i(t)\neq a_i(t-1)\}.
 \label{eq:foundation-switch}
\end{equation}
低切换率可以表示行为较稳定，但不能单独说明合作水平较高。全部节点长期背叛也可能具有较低切换率，因此需要与合作率共同解释。

对于参考对象选择，本文分别统计全部选择中的异常来源比例，以及真正触发正奖励优势修正的事件中的异常来源比例。前者反映候选排序，后者反映已进入更新的消息来源；二者分母不同。门控平均值表示连续缩放程度，实际附加更新的平均绝对值则同时受到门控、奖励优势和归一化分母影响。不能用门控值直接替代真实更新幅度。

\subsection{独立重复与对照设计}

基本重复单位为一次完整的多智能体仿真。同一网格中的节点通过重叠群组交互，同一次运行的时间步也具有相关性，不能将它们视为额外的独立样本。方法比较使用相同初始化和环境随机数构造配对运行，并分别记录每次运行的结果。若模型训练本身具有随机性，还需区分优化初始化产生的差异与固定模型在不同环境种子上的差异。

开发集用于实现检查、参数选择和方法调整，验证集用于比较候选模型，最终测试集用于冻结方法后的评价。开发结果可以指导下一步设计，但不能在反复查看和调整后仍被称为独立测试证据。对于存在不同长期状态的实验，还应展示逐次结果及其分布，避免一个平均值掩盖部分运行成功、部分运行失败的情况。

\section{本章小结}

本章分别定义了空间博弈、个体价值学习和邻居消息三个层次，并说明合作率、福利、切换率以及消息使用指标的含义。后续研究在相同问题背景下改变社会信息的使用方式：工作1分析可靠观测条件下的固定邻居影响规则，工作2研究受扰消息下的参考对象选择与修正控制。两项工作都需要保持动作、收益、更新时序和评价窗口的一致性，并在其实际信息边界内解释实验结论。
